\section*{Capítulo \ref{ch:g1:measure} --- Medir y medir el tiempo}

\begin{solution}{exo:g1:measure:1}
Depende de los lápices; al alinearlos por un extremo, el más largo llega
más lejos y el más corto se queda más cerca.
\end{solution}

\begin{solution}{exo:g1:measure:2}
Las respuestas varían; cada medida se lee poniendo el $0$ de la regla al principio
del objeto y leyendo el centímetro entero donde termina.
\end{solution}

\begin{solution}{exo:g1:measure:3}
Después del lunes: martes. Antes del domingo: sábado. Dos días después del
viernes: domingo.
\end{solution}

\begin{solution}{exo:g1:measure:4}
Ayer: miércoles. Mañana: viernes. Dentro de tres días: domingo.
\end{solution}

\begin{solution}{exo:g1:measure:5}
A las $6$ en punto la aguja corta apunta al $6$ y la aguja larga
apunta arriba, al $12$.
\end{solution}

\begin{solution}{exo:g1:measure:6}
Aguja corta en el $5$: las $5$ en punto. Aguja corta en el $11$: las $11$ en punto.
\end{solution}

\begin{solution}{exo:g1:measure:7}
Para $8$, por ejemplo $5 + 2 + 1$ o $2 + 2 + 2 + 2$ (hay otras respuestas
correctas).
\end{solution}

\begin{solution}{exo:g1:measure:8}
$5 + 2 + 2 + 2 = 11$. Leo tiene $11$.
\end{solution}

\begin{solution}{exo:g1:measure:9}
Del $3$ al $5$ hay $2$: a Zoe le devuelven $2$.
\end{solution}
